\documentclass{article}
\usepackage[T1]{fontenc}
\usepackage{lmodern}
\usepackage{graphicx}
\usepackage{amsmath,amssymb}
\usepackage[a4paper,margin=2cm]{geometry}
\usepackage[absolute,overlay]{textpos}
\setlength{\TPHorizModule}{1mm}
\setlength{\TPVertModule}{1mm}
\usepackage[indonesian]{babel}
\usepackage{hyperref}
\setlength{\emergencystretch}{2em}
% unit: Halaman 4

\begin{document}

% === Halaman 4 (python/native) ===

• Ukuran blok pesan = 64 bit

• Panjang kunci = 256 bit

• Jumlah putaran = 32 putaran

• Setiap putaran menggunakan kunci putaran. 

• Kunci putaran sebenarnya hanya ada 8 buah, K1 sampai K8,  namun  karena ada 

32 putaran, maka 8 buah kunci putaran ini dijadwalkan penggunaannya. 

• Standard GOST yang baru menspesifikasikan cipher blok 128-bit yang diberi 

nama Kuznyechik. 

4

Tabel 6.2  Penjadwalan kunci internal GOST 


Putaran 

1 

2 

3 

4 

5 

6 

7 

8 

9 

10 

11 

12 

13 

14 

15 

16 

Kunci internal 

K1 

K2 K3 K4 K5 K6 K7 K8 K1 K2 K3 K4 K5 K6 K7 K8 


Putaran 

17 

18 

19 

20 

21 

22 

23 

24 

25 

26 

27 

28 

29 

30 

31 

32 

Kunci internal 

K1 

K2 K3 K4 K5 K6 K7 K8 K8 K7 K6 K5 K4 K3 K2 K1

\end{document}
